\chapter{Cenni su Ottimizzazione e Scalabilità}

Nei capitoli precedenti abbiamo trattato il database relazionale come un'entità logica in grado di restituire i dati interrogati. Tuttavia, nell'ingegneria del software su larga scala, questa astrazione non è più sufficiente. Quando le tabelle passano da poche migliaia a decine di milioni di righe, un'interrogazione scritta in modo superficiale può causare la saturazione delle risorse e il blocco dell'intero sistema.

Questo capitolo si discosta dalla sintassi SQL per esplorare l'architettura fisica dello storage. Comprendere come il Database Management System (DBMS) organizza, indicizza e recupera i byte sul disco rigido o sull'SSD è il requisito fondamentale per poter progettare servizi backend realmente scalabili e per affrontare con successo i colloqui tecnici di livello avanzato.

\section{L'Anatomia dello Storage: Perché le query diventano lente?}

Per comprendere l'assoluta necessità delle complesse strutture dati utilizzate dai database relazionali moderni, è didatticamente utile partire dall'analisi della forma più elementare possibile di memorizzazione dei dati.

\paragraph{Il Database Primitivo: Il Log Append-Only}

Immaginiamo di voler costruire il nostro DBMS personale da zero. La struttura di memorizzazione più semplice in assoluto è un normale file di testo in cui ogni riga rappresenta un record (es. un file testuale separato da virgole). 

Quando un'applicazione backend invia un comando di \texttt{INSERT}, il nostro database rudimentale si limita ad aprire il file e ad accodare la nuova stringa di testo alla fine (operazione di \textit{append}).

\begin{observation}[Efficienza estrema in Scrittura]
L'accodamento sequenziale (\textit{Append-Only Log}) è in assoluto l'operazione di scrittura più veloce che un supporto di memoria fisico possa eseguire. Non richiedendo alcuna riorganizzazione dei dati preesistenti né complessi posizionamenti della testina del disco, il costo computazionale di un inserimento è minimo e rigorosamente costante nel tempo, denotato in notazione asintotica come $O(1)$.
\end{observation}

Il collasso architetturale di questo modello primitivo si manifesta inesorabilmente in fase di lettura. Quando l'applicazione invia una query di estrazione (es. \texttt{SELECT * FROM log WHERE id = 12345}), il database primitivo non ha alcuna informazione su dove si trovi fisicamente quel record. 

L'unica strategia algoritmica possibile per il motore è aprire il file dalla prima riga e scandirlo integralmente verso il basso finché non trova una corrispondenza. Questa operazione prende il nome di \textbf{Ricerca Lineare} (\textit{Sequential Scan}). Se il file è cresciuto fino a contenere 10 milioni di righe, il database potrebbe dover leggere 10 milioni di stringhe per trovare l'unica tupla richiesta. L'algoritmo ha un costo computazionale proporzionale alla mole dei dati, ovvero $O(N)$. All'aumentare del database, il tempo di risposta della singola query degrada in modo lineare, portando in breve tempo il backend al collasso prestazionale (il cosiddetto \textit{Timeout}).

\paragraph{Il Compromesso Fondamentale del Software Engineering}

Per evitare il disastroso costo della ricerca lineare $O(N)$, l'ingegneria dei dati adotta la stessa identica soluzione utilizzata nell'editoria stampata per trovare un concetto in un libro di mille pagine: si crea un \textbf{Indice}. 

Un indice, all'interno di un database, è una struttura dati aggiuntiva e fisicamente separata dai dati primari. Il suo unico scopo è mantenere una mappa ordinata (una "segnaletica") che indichi al motore di ricerca l'esatto indirizzo di memoria (\textit{offset}) di ciascun record. Invece di scansionare i dati, il DBMS interroga l'indice, trova le coordinate del dato, e lo recupera a colpo sicuro.

Sebbene l'idea appaia come la soluzione definitiva, la sua implementazione introduce nel sistema quello che Martin Kleppmann definisce il "compromesso fondamentale" (\textit{Trade-off}) dei database.Mantenere una struttura dati aggiuntiva ha un costo architetturale severo. Sebbene un indice abbatta drasticamente le tempistiche di lettura, esso \textbf{degrada le prestazioni in scrittura}. 

Ogni volta che l'applicazione esegue un'istruzione \texttt{INSERT}, \texttt{UPDATE} o \texttt{DELETE}, il DBMS non può più limitarsi ad accodare velocemente il dato nel file primario ($O(1)$), ma è fisicamente costretto a ri-calcolare, ri-bilanciare e aggiornare simultaneamente tutti gli indici associati a quella tabella affinché la mappa rimanga coerente.

Non esiste un database relazionale "ottimizzato per impostazione predefinita". La creazione degli indici è un'operazione manuale e rappresenta una precisa responsabilità dell'ingegnere software o del Database Administrator (DBA), il quale deve bilanciare l'infrastruttura basandosi sui reali \textit{pattern} di interrogazione dell'applicazione.

La regola aurea della scalabilità afferma che: \textit{"Un indice ben scelto velocizza enormemente le query di lettura, ma ogni singolo indice aggiunto rallenta proporzionalmente le operazioni di scrittura"}. Creare indici a tappeto su ogni colonna "per sicurezza" è un grave anti-pattern che rischia di paralizzare l'inserimento dei dati nel sistema.

\section{Strutture Dati per l'Indicizzazione}

Per evitare la scansione lineare, il DBMS deve organizzare gli indici tramite strutture dati avanzate. Sebbene esistano diversi approcci (come i log LSM utilizzati nei DB NoSQL), il mondo relazionale è dominato da decenni da un unico, indiscusso standard algoritmico.

\subsection{Il motore universale relazionale: Il B-Tree (Balanced Tree)}

Introdotto nei primi anni '70 e costantemente raffinato, il \textbf{B-Tree} (precisamente nella sua variante $B^+$-Tree) è l'albero di ricerca bilanciato su cui poggia l'intera architettura di PostgreSQL. 

A differenza degli alberi binari di ricerca studiati nei corsi di informatica teorica, il B-Tree è progettato specificamente per lavorare in simbiosi con l'hardware dei supporti di memoria fisici (Dischi Meccanici ed SSD).

\paragraph{1. L'organizzazione fisica sul disco: Le Pagine}

Il segreto del B-Tree risiede nella sua granularità. Mentre i log di testo crescono in modo indefinito, il B-Tree scompone l'intero database in blocchi di memoria a dimensione fissa, chiamati \textbf{Pagine} (o \textit{Blocks}), tipicamente di $4\text{ KB}$ o $8\text{ KB}$ (in PostgreSQL il default è $8\text{ KB}$). Il DBMS legge o scrive sul disco sempre una pagina intera alla volta.

La struttura si sviluppa in modo gerarchico e piramidale:

\begin{itemize}
    \item \textbf{Nodo Radice (Root Node):} È la pagina iniziale dell'indice. Contiene una sequenza ordinata di chiavi e una serie di puntatori (indirizzi di memoria del disco) che rimandano alle pagine del livello inferiore.
    \item \textbf{Nodi Interni (Internal Nodes):} Pagine intermedie che fungono da "svincoli autostradali". Leggendo la chiave cercata, indirizzano il motore verso la pagina successiva, restringendo progressivamente il campo di ricerca.
    \item \textbf{Nodi Foglia (Leaf Nodes):} Sono le pagine situate alla base della piramide. Contengono le chiavi effettive e, accanto a ciascuna di esse, il puntatore fisico (\textit{Tuple ID} o \textit{offset}) che indica dove si trova la riga reale all'interno della tabella principale. Nelle varianti moderne, tutte le foglie sono collegate tra loro da puntatori orizzontali per velocizzare la scansione sequenziale ordinata.
\end{itemize}



Il numero di puntatori a figli che una singola pagina può ospitare prende il nome di \textbf{Branching Factor} (o \textit{Fan-out}). Poiché una pagina di $8\text{ KB}$ può contenere centinaia di chiavi numeriche, il fattore di ramificazione è straordinariamente elevato (spesso superiore a 200–300).

\begin{example}[Tracciamento di una ricerca puntuale nel B-Tree]
Supponiamo che il backend Java esegua l'interrogazione \texttt{SELECT * FROM sinistro WHERE id\_cliente = 83}. La colonna \texttt{id\_cliente} è indicizzata tramite un B-Tree. Vediamo come il motore di PostgreSQL naviga fisicamente le pagine da $8\text{ KB}$ memorizzate sul disco per recuperare il record.

\textbf{Fase 1: Ispezione del Nodo Radice (Pagina \#100)}
Il motore carica in RAM la pagina di radice. Al suo interno trova una guida sintetica suddivisa in intervalli di chiavi:
\begin{itemize}
    \item Chiavi $< 50 \longrightarrow$ Punta alla Pagina \#101
    \item Chiavi da $50$ a $99 \longrightarrow$ Punta alla Pagina \#102
    \item Chiavi $\ge 100 \longrightarrow$ Punta alla Pagina \#103
\end{itemize}
Poiché il valore cercato è $83$, il Query Planner scarta immediatamente le pagine \#101 e \#103 (e tutti i sotto-alberi ad esse collegati). Il motore esegue una lettura sul disco per caricare solo la \textbf{Pagina \#102}.

\textbf{Fase 2: Navigazione del Nodo Interno (Pagina \#102)}
La Pagina \#102 è un nodo intermedio (lo svincolo). Al suo interno la granularità aumenta e l'intervallo viene ulteriormente frammentato:
\begin{itemize}
    \item Chiavi da $50$ a $69 \longrightarrow$ Punta alla Pagina \#201
    \item Chiavi da $70$ a $89 \longrightarrow$ Punta alla Pagina \#202
    \item Chiavi da $90$ a $99 \longrightarrow$ Punta alla Pagina \#203
\end{itemize}
Il valore $83$ ricade nell'intervallo centrale. Il motore ignora le pagine \#201 e \#203 ed effettua un secondo accesso al disco per caricare la \textbf{Pagina \#202}.

\textbf{Fase 3: Identificazione nel Nodo Foglia (Pagina \#202)}
La Pagina \#202 è un nodo foglia, la base della piramide. Qui non ci sono più puntatori ad altre pagine dell'indice, ma le chiavi reali associate agli indirizzi fisici dei dati nella tabella (\textit{Heap File}):
\begin{itemize}
    \item Chiave $75 \longrightarrow$ Record memorizzato all'offset \texttt{0x0A2}
    \item Chiave $83 \longrightarrow$ Record memorizzato all'offset \texttt{0x1F4}
    \item Chiave $89 \longrightarrow$ Record memorizzato all'offset \texttt{0x33B}
\end{itemize}
Il motore trova la corrispondenza esatta per la chiave $83$ e ottiene il puntatore fisico \texttt{0x1F4}.

\textbf{Fase 4: Estrazione dal File dei Dati (Heap File)}
L'indice ha terminato il suo compito. Il database esegue l'ultimo accesso al disco puntando direttamente all'offset \texttt{0x1F4} della tabella \texttt{sinistro}, estrae la riga completa (targa, data, danno stimato) e la impacchetta nel \textit{Result Set} da inviare sulla rete all'applicazione Java.
\end{example}


Grazie a un elevato \textit{Branching Factor}, l'albero si sviluppa pochissimo in altezza (profondità) ma immensamente in larghezza. Un B-Tree a soli 4 livelli logici, con un fattore di ramificazione pari a 250, è matematicamente in grado di ospitare oltre 3.9 miliardi di record. 

Cercare un ID specifico in un database da 4 miliardi di righe non richiede 4 miliardi di scansioni ($O(N)$), ma l'apertura di \textbf{sole 4 pagine sul disco}. Il costo computazionale è rigidamente logaritmico, denotato come $O(\log N)$. Le prestazioni rimangono stabili, fulminee e indipendenti dalla crescita del database.


\paragraph{2. Il comportamento in Scrittura: Il Page Split}

Il B-Tree viene definito "Bilanciato" perché l'algoritmo garantisce che tutte le pagine foglia si trovino sempre, in ogni istante di tempo, alla stessa identica distanza (profondità) dalla radice. Questo previene il degrado dell'albero.

Tuttavia, mantenere questo perfetto bilanciamento matematico durante un'operazione di \texttt{INSERT} richiede un calcolo computazionale oneroso. Quando l'applicazione backend inserisce una nuova riga, il DBMS naviga il B-Tree per individuare la pagina foglia corretta in cui allocare la nuova chiave in ordine numerico. 

A questo punto si aprono due scenari fisici:

\begin{enumerate}
    \item \textbf{Spazio disponibile:} Se la pagina da $8\text{ KB}$ ha ancora byte liberi, la chiave viene semplicemente scritta nella pagina. Il costo è minimo.
    \item \textbf{Pagina satura (Il Page Split):} Se la pagina è fisicamente piena, il motore non può semplicemente accodare il dato. Deve attivare l'algoritmo di \textbf{Page Split} (Scissione della Pagina): la pagina viene allocata \textit{ex novo} sul disco, i dati vengono divisi a metà tra la vecchia e la nuova pagina, e la chiave mediana viene spinta verso l'alto nel nodo padre per aggiornare i puntatori.
\end{enumerate}

\begin{warningbox}[Il collasso da Page Split a cascata]
Se anche il nodo padre è saturo, il fenomeno del \textit{Page Split} si propaga ricorsivamente verso l'alto, potendo arrivare a sdoppiare il nodo radice stesso e aumentando l'altezza dell'intero albero. 

Durante un inserimento che causa un Page Split a cascata, il database deve effettuare molteplici operazioni di scrittura in punti diversi del disco rigido anziché un singolo accodamento rapido. Questo spiega sperimentalmente perché le \texttt{INSERT} massive rallentino drasticamente in presenza di troppi indici attivi.
\end{warningbox}
